\chapter{Compactness of Riemannian Manifolds and Flows}

\section{Convergence and compactness of manifolds}

It is reasonable to suggest that a sequence $\{g_i\}$ of Riemannian metrics on
a manifold $\M$ should converge to a metric $g$ when $g_i \to g$ as tensors.
However, we would like a notion of convergence of Riemannian manifolds which is
diffeomorphism invariant: it should not be affected if we modify each metric
$g_i$ by an $i$-dependent diffeomorphism. Once we have asked for such
invariance, it is necessary to discuss convergence with respect to a point of
reference on each manifold.

\begin{definition}
\label{topping-ch7-pointed-cheeger-gromov-convergence}
\dcref{7.1.1}
\group{compactness-manifolds}
\source{topping:page-0080}
A sequence $(\M_i,g_i,p_i)$ of smooth, complete, pointed Riemannian manifolds
(that is, Riemannian manifolds $(\M_i,g_i)$ and points $p_i \in \M_i$) is said
to \emph{converge} (smoothly) to the smooth, complete, pointed manifold
$(\M,g,p)$ as $i \to \infty$ if there exist
\begin{enumerate}
\item[(i)] a sequence of compact sets $\Omega_i \subset \M$, exhausting
$\M$ (that is, so that any compact set $K \subset \M$ satisfies
$K \subset \Omega_i$ for sufficiently large $i$) with $p \in \operatorname{int}(\Omega_i)$
for each $i$;
\item[(ii)] a sequence of smooth maps $\phi_i : \Omega_i \to \M_i$ which are
diffeomorphic onto their image and satisfy $\phi_i(p) = p_i$ for all $i$;
\end{enumerate}
such that
\[
\phi_i^* g_i \to g
\]
smoothly as $i \to \infty$ in the sense that for all compact sets
$K \subset \M$, the tensor $\phi_i^* g_i - g$ and its covariant derivatives of
all orders (with respect to any fixed background connection) each converge
uniformly to zero on $K$.
\end{definition}

To demonstrate why we need to include the points $p_i$ in the above
definition, consider the following example.

\begin{remark}
\label{topping-ch7-pointed-convergence-example}
\dcref{7.1.2}
\group{compactness-manifolds}
\source{topping:page-0081}
Suppose first that for every $i$, $(\M_i,g_i)$ is equal to the
$(\mathcal{N},h)$ as shown above. Then $(\M_i,g_i,q) \to (\mathcal{N},h,q)$,
but $(\M_i,g_i,s_i)$ converges
to a cylinder.
\end{remark}

\begin{remark}
\label{topping-ch7-compact-limit-possible}
\dcref{7.1.2}
\group{compactness-manifolds}
\source{topping:page-0081}
It is possible to have $\M_i$ compact for all $i$, but have the limit $\M$
non-compact.
\end{remark}

Two consequences of the convergence $(\M_i,g_i,p_i) \to (\M,g,p)$ are that
\begin{enumerate}
\item[(i)] for all $s > 0$ and $k \in \{0\} \cup \NN$,
\[
\sup_{i \in \NN} \sup_{B_{g_i}(p_i,s)} \lvert \nabla^k \Rm(g_i)\rvert < \infty;
\tag{7.1.1}
\]
\item[(ii)]
\[
\inf_i \inj(\M_i,g_i,p_i) > 0,
\tag{7.1.2}
\]
\end{enumerate}
where $\inj(\M_i,g_i,p_i)$ denotes the injectivity radius of
$(\M_i,g_i)$ at $p_i$.

In fact, conditions (i) and (ii) are sufficient for subconvergence. Various
incarnations of the following result appear in, or can be derived from, papers
of (for example) Greene and Wu [16], Fukaya [12] and Hamilton [22], all of
which can be traced back to original ideas of Gromov [17] and Cheeger.

\begin{theorem}
\label{topping-ch7-compactness-manifolds}
\dcref{7.1.3}
\group{compactness-manifolds}
\source{topping:page-0081}
\uses{topping-ch7-pointed-cheeger-gromov-convergence,topping-ch7-pointed-convergence-example,topping-ch7-compact-limit-possible}
Suppose that $(\M_i,g_i,p_i)$ is a sequence of complete, smooth, pointed
Riemannian manifolds (all of dimension $n$) satisfying (7.1.1) and (7.1.2).
Then there exists a complete, smooth, pointed Riemannian manifold
$(\M,g,p)$ (of dimension $n$) such that after passing to some subsequence in
$i$,
\[
(\M_i,g_i,p_i) \to (\M,g,p).
\]
\end{theorem}

\begin{remark}
\label{topping-ch7-curvature-control-needed}
\dcref{7.1.4}
\group{compactness-manifolds}
\source{topping:page-0082}
Most theorems of this type only require boundedness of $\lvert \Rm\rvert$
(or even $\lvert \Ric\rvert$) and then ask for weaker convergence. Our
manifolds will always arise as time-slices of Ricci flows and so we may as
well assume $\lvert \nabla^k \Rm\rvert$ bounds by our derivative estimates
from Theorem~3.3.1.
\end{remark}

There are other notions of convergence which can handle this type of limit. For
example, Gromov-Hausdorff convergence allows us to take limits of metric
spaces. (See, for example [3].)

\begin{remark}
\label{topping-ch7-injectivity-radius-needed}
\dcref{7.1.5}
\group{compactness-manifolds}
\source{topping:page-0082}
The uniform positive lower bound on the injectivity radius is also necessary,
since otherwise we could have, for example, degeneration of cylinders.
\end{remark}

\begin{remark}
\label{topping-ch7-injectivity-radius-propagation}
\dcref{7.1.6}
\group{compactness-manifolds}
\source{topping:page-0082}
Given curvature bounds as in (7.1.1), the injectivity radius lower bound at
$p_i$ implies a positive lower bound for the injectivity radius at other
points $q \in \M_i$ in terms of the distance from $p_i$ to $q$, and the
curvature bounds, as discussed in [22], say.
\end{remark}

\begin{remark}
\label{topping-ch7-compactness-usage}
\dcref{7.1.7}
\group{compactness-manifolds}
\source{topping:page-0082}
In these notes, whenever we wish to apply Theorem~7.1.3, we will know not only
(7.1.1) but also that for all $k \in \{0\} \cup \NN$,
\[
\sup_i \sup_{x \in \M} \lvert \nabla^k \Rm(g_i)\rvert < \infty.
\]
\end{remark}

\section{Convergence and compactness of flows}

One can derive, from the compactness theorem for manifolds (Theorem~7.1.3) a
compactness theorem for Ricci flows.

\begin{definition}
\label{topping-ch7-flow-convergence}
\dcref{7.2.1}
\group{compactness-flows}
\source{topping:page-0083}
Let $(\M_i,g_i(t))$ be a sequence of smooth families of complete Riemannian
manifolds for $t \in (a,b)$ where $-\infty \le a < 0 < b \le \infty$. Let
$p_i \in \M_i$ for each $i$. Let $(\M,g(t))$ be a smooth family of complete
Riemannian manifolds for $t \in (a,b)$ and let $p \in \M$. We say that
\[
(\M_i,g_i(t),p_i) \to (\M,g(t),p)
\]
as $i \to \infty$ if there exist
\begin{enumerate}
\item[(i)] a sequence of compact $\Omega_i \subset \M$ exhausting $\M$ and
satisfying $p \in \operatorname{int}(\Omega_i)$ for each $i$;
\item[(ii)] a sequence of smooth maps $\phi_i : \Omega_i \to \M_i$, diffeomorphic
onto their image, and with $\phi_i(p)=p_i$;
\end{enumerate}
such that
\[
\phi_i^{*} g_i(t) \to g(t)
\]
as $i \to \infty$ in the sense that $\phi_i^{*} g_i(t)-g(t)$ and its
derivatives of every order (with respect to time as well as covariant space
derivatives with respect to any fixed background connection) converge
uniformly to zero on every compact subset of $\M \times (a,b)$.
\end{definition}

\begin{remark}
\label{topping-ch7-flow-backward-interval}
\dcref{7.2.2}
\group{compactness-flows}
\source{topping:page-0083}
It also makes sense to talk about convergence on (for example) the time
interval $(-\infty,0)$, even when flows are defined only for $(-T_i,0)$ with
$T_i \to \infty$.
\end{remark}

Hamilton then proves the following result [22], based on the core
Theorem~7.1.3.

\begin{theorem}
\label{topping-ch7-compactness-ricci-flows}
\dcref{7.2.3}
\group{compactness-flows}
\source{topping:page-0083,topping:page-0084}
Suppose that $\M_i$ is a sequence of manifolds of dimension $n$, and let
$p_i \in \M_i$ for each $i$. Suppose that $g_i(t)$ is a sequence of complete
Ricci flows on $\M_i$ for $t \in (a,b)$, where $-\infty \le a < 0 < b \le
\infty$. Suppose that
\begin{enumerate}
\item[(i)] $\sup_i \sup_{x \in \M_i,\, t \in (a,b)} \lvert \Rm(g_i(t))(x)\rvert < \infty$; and
\item[(ii)] $\inf_i \inj(\M_i,g_i(0),p_i) > 0$.
\end{enumerate}
Then there exist a manifold $\M$ of dimension $n$, a complete Ricci flow
$g(t)$ on $\M$ for $t \in (a,b)$, and a point $p \in \M$ such that, after
passing to a subsequence in $i$,
\[
(\M_i,g_i(t),p_i) \to (\M,g(t),p).
\]
\end{theorem}

\section{Blowing up at singularities I}

As we have mentioned before, one application of the compactness of Ricci flows
in Theorem~7.2.3 that we have in mind, is to analyse rescalings of Ricci flows
near their singularities. Let $(\M,g(t))$ be a Ricci flow with $\M$ closed, on
the maximal time interval $[0,T)$ -- as defined at the end of
Section~5.2 -- with $T < \infty$. By Theorem~5.3.1 we know that
\[
\sup_{\M} \lvert \Rm\rvert(\cdot,t) \to \infty
\]
as $t \uparrow T$. Let us choose points $p_i \in \M$ and times $t_i \uparrow T$
such that
\[
\lvert \Rm\rvert(p_i,t_i) = \sup_{x \in \M,\, t \in [0,t_i]} \lvert \Rm\rvert(x,t),
\]
by, for example, picking $(p_i,t_i)$ to maximise $\lvert \Rm\rvert$ over
$\M \times [0, T - \frac{1}{i}]$. Notice in particular that
$\lvert \Rm\rvert(p_i,t_i) \to \infty$ as $i \to \infty$. Define rescaled (and
translated) flows $g_i(t)$ by
\[
g_i(t) = \lvert \Rm\rvert(p_i,t_i)\, g\!\left(t_i + \frac{t}{\lvert \Rm\rvert(p_i,t_i)}\right).
\]
By the discussion in Section~1.2.3, $(\M,g_i(t))$ is a Ricci flow on the time
interval
\[
[-t_i\lvert \Rm\rvert(p_i,t_i), (T-t_i)\lvert \Rm\rvert(p_i,t_i)).
\]

Moreover, for each $i$, $\lvert \Rm(g_i(0))\rvert(p_i) = 1$ and for $t \le 0$,
\[
\sup_{\M} \lvert \Rm(g_i(t))\rvert \le 1.
\]
Furthermore, by Theorem~3.2.11, for $t > 0$,
\[
\sup_{\M} \lvert \Rm(g_i(t))\rvert \le \frac{1}{1 - C(n)t}.
\]
Therefore, for all $a < 0$ and some $b = b(n) > 0$, $g_i(t)$ is defined for
$t \in (a,b)$ and
\[
\sup_i \ \sup_{\M \times (a,b)} \lvert \Rm(g_i(t))\rvert < \infty.
\]

By Theorem~7.2.3, we can pass to a subsequence in $i$, and get convergence
\[
(\M, g_i(t), p_i)\to (\mathcal{N}, \hat g(t), p_\infty)
\]
to a ``singularity model'' Ricci flow $(\mathcal{N},\hat g(t))$, provided that
we can
establish the injectivity radius estimate
\begin{equation}
\inf_i \inj(\M, g_i(0), p_i) > 0.
\tag{7.3.1}
\end{equation}

\noindent\textbf{Example 7.3.1.} As an example of what we have in mind, recall
the possibility of a `neck pinch' as discussed in Section~1.3.2. The limit of
the rescaled flows should, in principle, be a shrinking cylinder.

\noindent
The missing step for us at this stage is to get the lower bound on the
injectivity radius stated in (7.3.1). Historically, this has been a major
difficulty, except in some special cases. However, as we will see in the next
chapter, this issue has been resolved recently by Perelman.
